% =====================================================================
%  CAPÍTULO 3 — METODOLOGÍA
% =====================================================================
\chapter{Metodología}
\label{cap:metodologia}

La naturaleza de este trabajo —que aplica un instrumental analítico poco habitual en un TFG de
Administración y Dirección de Empresas— aconseja dedicar un capítulo completo a la
metodología, tal como recomienda la propia guía de la Facultad para los trabajos cuyo método
exige ser explicado con detalle. En las páginas que siguen se describe, de forma ordenada, el
diseño de la investigación, la manera en que se construyeron la muestra y el corpus documental,
las técnicas de procesamiento del lenguaje natural empleadas —con cita expresa de la fuente
original de cada una—, la construcción del índice de \textit{greenwashing} y, por último, el
aparato estadístico utilizado para contrastar las hipótesis. El objetivo es que el lector
comprenda no solo \textit{qué} se ha hecho, sino \textit{por qué} se ha hecho de ese modo y qué
se persigue con cada decisión.

\section{Diseño de la investigación y preguntas}
\label{sec:diseno}

El trabajo se inscribe en la tradición del análisis de contenido, pero lo aborda de manera
cuantitativa y a gran escala mediante técnicas automáticas. Su unidad de análisis es la
subsección de sostenibilidad de cada informe corporativo, y su diseño combina dos dimensiones.
La dimensión transversal permite comparar empresas de distintos sectores y países en un mismo
momento, y da respuesta a las preguntas sobre diferencias y determinantes (RQ2 y RQ3). La
dimensión temporal —tres ejercicios consecutivos, 2022, 2023 y 2024— permite observar la
evolución del reporting en el tránsito de la NFRD a la CSRD y sostiene la cuarta pregunta
(RQ4). El conjunto conforma un panel de empresa por año sobre el que se proyectan todas las
métricas textuales.

Las cuatro preguntas de investigación, ya enunciadas en la introducción, se corresponden con
los cuatro objetivos específicos y guían la organización del capítulo de resultados: los temas
predominantes (RQ1), las diferencias por sector y país (RQ2), los determinantes de la
especificidad y la posible relación entre tono y concreción (RQ3) y la evolución entre
regímenes normativos (RQ4). Conviene insistir en la decisión metodológica que vertebra todo el
trabajo: el \textit{greenwashing} no se mide mediante calificaciones ESG de terceros, sino que
se operacionaliza a partir de variables textuales internas —tono, especificidad, vaguedad,
cobertura cuantitativa— extraídas del propio informe. Esta opción preserva la independencia
entre el objeto analizado y el instrumento de medida y evita la circularidad en que incurriría
cualquier análisis que empleara una nota externa de sostenibilidad como medida de la calidad de
la comunicación.

\section{Muestra: el STOXX Europe 600}
\label{sec:muestra}

El universo de estudio es el STOXX Europe 600, índice que reúne a seiscientas de las mayores
empresas cotizadas de Europa y que constituye una referencia habitual del mercado europeo de
renta variable \parencite{stoxx2024}. De ese universo se seleccionó una muestra de
\textbf{196 empresas} mediante un muestreo estratificado por sector, con una semilla aleatoria
fija para garantizar la reproducibilidad. La estratificación responde a una razón de fondo: el
índice está fuertemente concentrado en unos pocos sectores —servicios financieros e
industriales suman por sí solos más del cuarenta por ciento—, de modo que un muestreo aleatorio
simple sobrerrepresentaría a estos y dejaría sin apenas presencia a sectores pequeños pero muy
relevantes para el análisis de sostenibilidad, como las utilities o el sector inmobiliario.
Fijando una cuota por sector se garantiza que todos aporten al análisis comparativo. Dentro de
cada estrato la selección fue aleatoria, y no por capitalización, para evitar el sesgo hacia
las mayores compañías —dotadas de más recursos y bajo mayor presión mediática— y captar así la
verdadera dispersión en la calidad del reporting. Se aplicó, además, un límite geográfico que
corrige la sobrerrepresentación estructural del Reino Unido en el índice.

La muestra resultante distribuye sus 196 empresas entre once supersectores y diecisiete países,
tal como recogen las tablas~\ref{tab:muestra-sector} y~\ref{tab:muestra-pais}. Cada empresa se
observa en los tres ejercicios, lo que da lugar a un panel de 588 observaciones de empresa por
año. Para cada una de ellas se recopilaron, además del informe corporativo, los datos
financieros básicos —capitalización, rentabilidad sobre activos y sobre fondos propios,
endeudamiento, cifra de negocio— procedentes de Yahoo Finance, que se emplean como variables de
control en el análisis de los determinantes \parencite{yfinance2024}.

\begin{table}[H]
\centering
\caption{Distribución de la muestra por supersector}
\label{tab:muestra-sector}
\begin{tabular}{lr@{\hspace{2cm}}lr}
\toprule
\textbf{Supersector} & \textbf{Empresas} & \textbf{Supersector} & \textbf{Empresas} \\
\midrule
Consumer Discretionary & 39 & Consumer Staples & 17 \\
Financials & 33 & Technology & 11 \\
Industrials & 20 & Energy & 10 \\
Basic Materials & 18 & Health Care & 10 \\
Communication Services & 18 & Real Estate & 10 \\
 & & Utilities & 10 \\
\bottomrule
\end{tabular}
\\[4pt]
{\footnotesize Fuente: elaboración propia. Total: 196 empresas.}
\end{table}

\begin{table}[H]
\centering
\caption{Distribución de la muestra por país}
\label{tab:muestra-pais}
\begin{tabular}{lr@{\hspace{1.2cm}}lr@{\hspace{1.2cm}}lr}
\toprule
\textbf{País} & \textbf{Emp.} & \textbf{País} & \textbf{Emp.} & \textbf{País} & \textbf{Emp.} \\
\midrule
Reino Unido & 30 & Italia & 14 & Austria & 4 \\
Alemania & 29 & Suecia & 11 & Dinamarca & 3 \\
Francia & 28 & Noruega & 8 & Irlanda & 2 \\
Suiza & 19 & Bélgica & 7 & Portugal & 2 \\
España & 17 & Finlandia & 5 & Luxemburgo & 1 \\
Países Bajos & 15 & & & Israel & 1 \\
\bottomrule
\end{tabular}
\\[4pt]
{\footnotesize Fuente: elaboración propia. Total: 196 empresas.}
\end{table}

La elección de los tres ejercicios no es arbitraria. El año 2022 representa el último ejercicio
mayoritariamente regido por la NFRD y funciona como línea de base anterior a la CSRD; 2023
recoge informes elaborados bajo expectativas crecientes de la nueva norma; y 2024 es el primer
ejercicio de reporte obligatorio bajo la CSRD para las grandes empresas de interés público.
La comparación entre los extremos de esa ventana —2022 frente a 2024— es la que permite
responder a la pregunta sobre la evolución del reporting.

\section{Construcción del corpus}
\label{sec:corpus}

Disponer de una muestra de empresas no equivale a disponer de un material analizable: los
informes corporativos son documentos extensos, publicados en formato PDF, con estructuras
heterogéneas y, con frecuencia, en idiomas distintos del inglés. Convertir ese material en un
corpus de texto limpio y homogéneo fue una fase de trabajo en sí misma, que conviene describir
porque condiciona la validez de todo lo que sigue.

El primer paso fue la \textbf{extracción del texto} de los informes. La mayoría de los
documentos permite una extracción directa, pero algunos presentaban el texto como imagen o con
una codificación de caracteres defectuosa que lo hacía ilegible para las herramientas
automáticas. En esos casos se recurrió a un reconocimiento óptico de caracteres aplicado
únicamente a las páginas problemáticas, conservando el texto nativo en el resto, de modo que la
calidad de la fuente original se preservó siempre que fue posible. El segundo paso fue la
\textbf{homogeneización idiomática}. Dado que los modelos de lenguaje empleados más adelante
están entrenados sobre texto en inglés, los informes publicados en otros idiomas se
sustituyeron por la versión oficial en inglés del mismo emisor —habitualmente el
\textit{Annual Report} o el documento de registro universal—, nunca por una traducción
automática, que habría introducido un ruido incompatible con el análisis de sentimiento. El
corpus final quedó así íntegramente en inglés y libre de traducciones artificiales.

El tercer y más delicado paso fue la \textbf{segmentación}. Cada informe se dividió para aislar,
por un lado, el informe de gestión y, por otro, la subsección específicamente dedicada a la
sostenibilidad, que es la unidad de análisis de este trabajo. Esta separación se realizó
combinando dos criterios: el índice navegable del propio PDF, cuando estaba disponible, y, en su
defecto, la densidad de vocabulario de sostenibilidad por página. Una decisión importante de
este proceso merece subrayarse: como la subsección de sostenibilidad está contenida dentro del
informe de gestión, se optó por definir el informe de gestión \textit{excluyendo} la parte de
sostenibilidad, de manera que ambas secciones no solapen y se evite el doble cómputo del mismo
texto. A cada documento se le asignó, además, una etiqueta de fiabilidad de la segmentación,
que permite identificar los casos en que la subsección se aisló de forma parcial y someterlos a
un análisis de sensibilidad.

El resultado de todo este proceso es un corpus de \textbf{586 documentos} de sostenibilidad
analizables, correspondientes a las 196 empresas a lo largo de los tres ejercicios. La
diferencia respecto de las 588 observaciones del panel se debe a la exclusión de dos informes
que, pese a existir, carecían de contenido de sostenibilidad analizable. Como ilustra la
evolución del tamaño de los documentos —que se examina en el capítulo de resultados—, el
material de partida no es menor: la subsección de sostenibilidad alcanza una mediana de varios
miles de palabras por documento y crece de forma acusada a lo largo del periodo.

\section{Técnicas de procesamiento del lenguaje natural}
\label{sec:tecnicas}

Sobre el corpus así construido se aplicó un conjunto de técnicas complementarias. La lógica que
las articula es la triangulación: en lugar de confiar en un único método, se combinan varios de
naturaleza distinta —léxicos, estadísticos y basados en modelos de lenguaje— de modo que sus
resultados puedan contrastarse entre sí y reforzarse mutuamente. A continuación se describen
agrupadas según su finalidad.

\subsection{Diccionarios: tono y cobertura temática}

El primer bloque de técnicas se apoya en diccionarios, esto es, en listas de términos
previamente clasificados. Su gran virtud es la transparencia: el resultado es plenamente
auditable, pues siempre puede rastrearse qué palabras lo han generado. Se emplearon dos. El
primero es el diccionario financiero de \textcite{loughran2011}, que clasifica miles de
términos en categorías —negativo, positivo, incertidumbre, lenguaje litigioso, modalidad fuerte
y débil— en su versión más reciente \parencite{lmdict2024}, y que permite caracterizar el tono y,
sobre todo, el grado de cautela o vaguedad del discurso. El segundo es un
\textbf{diccionario propio de términos ESRS}, construido a partir de los requisitos de
divulgación de los doce estándares —tomando como referencia la taxonomía oficial publicada por
EFRAG \parencite{efrag2023}— y organizado en once categorías temáticas, con el que se mide la
amplitud con que cada informe cubre los distintos pilares del marco europeo. Este enfoque sigue
la estela de \textcite{suta2025}, que validan el uso de diccionarios para evaluar la cobertura de
los temas ESRS. Conviene precisar qué mide y qué no mide este
indicador: capta la \textit{anchura} del vocabulario empleado, no la profundidad ni el
cumplimiento normativo, de modo que su lectura debe ser siempre relativa y comparativa.

\subsection{Modelado de temas: ¿de qué hablan los informes?}

El segundo bloque busca descubrir, sin imponerlos de antemano, los asuntos que vertebran el
discurso de sostenibilidad. Para ello el corpus se troceó en párrafos y se aplicaron dos
técnicas de modelado de temas de naturaleza distinta, lo que permite contrastar sus resultados.
La primera, el \textit{Latent Dirichlet Allocation} o LDA \parencite{blei2003}, es un método
probabilístico clásico que representa cada documento como una mezcla de temas y cada tema como
una distribución de palabras; el número óptimo de temas se determinó mediante una medida de
coherencia que evalúa hasta qué punto las palabras de cada tema tienden a aparecer juntas en el
texto real. La segunda, BERTopic \parencite{grootendorst2022}, es un enfoque más reciente que,
en lugar de contar palabras, agrupa los párrafos según su \textit{significado}: para ello
representa cada fragmento de texto mediante un modelo de lenguaje \parencite{reimers2019},
reduce la complejidad de esa representación \parencite{mcinnes2018} y agrupa los fragmentos
semánticamente próximos \parencite{campello2013}. La ventaja de emplear ambos métodos es que,
si dos técnicas tan diferentes coinciden en identificar los mismos grandes temas, la confianza
en ese hallazgo aumenta.

\subsection{Modelos de lenguaje: sentimiento y especificidad}

El tercer bloque recurre a modelos de lenguaje especializados, herederos de la arquitectura
BERT \parencite{devlin2019,liu2019}, que clasifican cada frase del corpus según diversos
criterios. FinBERT \parencite{araci2019} evalúa el tono financiero general de cada frase
—positivo, negativo o neutro—. Las cuatro variantes de ClimateBERT
\parencite{webersinke2022}, las mismas que emplean Bingler y sus colegas, operan en cascada:
la primera detecta si una frase trata de clima; sobre las que sí lo hacen, las restantes
evalúan su sentimiento —riesgo frente a oportunidad—, su grado de compromiso —si contiene un
compromiso concreto o una mera declaración— y su especificidad —si aporta cifras, fechas y
metas o se queda en lo genérico—. Por último, un clasificador adicional, FinBERT-ESG-9
\parencite{yang2020}, asigna cada frase a una de nueve categorías temáticas de sostenibilidad,
lo que ofrece una tercera vía, independiente de los diccionarios y del modelado de temas, para
caracterizar el contenido del corpus. El uso de estos modelos sobre frases individuales —cientos
de miles en total— es lo que permite descender del nivel agregado del documento al detalle del
discurso.

\section{Construcción del índice de \textit{greenwashing}}
\label{sec:gwindex}

Las métricas anteriores capturan, cada una, una faceta del discurso de sostenibilidad. Para
sintetizar las señales de \textit{cheap talk} en una única medida comparable se construyó un
índice compuesto, denominado \gwindex{}, que combina cuatro de esos componentes tras
estandarizarlos. La definición operativa responde directamente a la noción de
\textit{greenwashing} adoptada en el trabajo y a las dimensiones de calidad de la divulgación de
\textcite{michelon2015}: un discurso es tanto más sospechoso de \textit{cheap talk} cuanto más
cauteloso y evasivo sea su lenguaje, cuantas más promesas de futuro formule sin respaldarlas con
cifras, cuanto menos contenido cuantitativo aporte y cuanto menos específica sea su información
climática. Formalmente, el índice se define como

\[
\textit{GW\_index} = z(\textit{hedging}) + z(\textit{promesa vaga})
- z(\textit{ratio cuantitativo}) - z(\textit{especificidad climática}),
\]

\noindent donde $z(\cdot)$ denota la estandarización de cada componente sobre el conjunto del
panel. El \textit{hedging} mide la proporción de lenguaje de incertidumbre y modalidad débil; la
\textit{promesa vaga} (\textit{ratio futuro sin cifra}), la proporción de frases prospectivas
que no contienen ningún dato numérico; el \textit{ratio cuantitativo}, la proporción de frases
que sí incorporan cifras concretas; y la \textit{especificidad climática}, la proporción de
frases climáticas que ClimateBERT clasifica como específicas. Los dos primeros componentes
suman al índice —más vaguedad, más sospecha— y los dos últimos restan —más concreción, menos
sospecha—. Un valor alto del \gwindex{}, por tanto, señala un reporting comparativamente más
cauteloso, más vago en sus promesas y menos respaldado por datos. El tono financiero general de
FinBERT se mantiene deliberadamente al margen del índice, para poder usarlo después como
variable independiente con la que contrastar la hipótesis del optimismo sin sustancia.

\section{Análisis estadístico inferencial}
\label{sec:estadistica}

El último componente metodológico es el aparato estadístico con el que se contrastan las
hipótesis. Para examinar las diferencias por sector y por país (RQ2) se empleó la prueba de
Kruskal-Wallis, un contraste no paramétrico apropiado cuando las distribuciones no son normales,
que evita imponer supuestos difíciles de sostener sobre variables como el \gwindex{}. Para
estudiar la evolución entre 2022 y 2024 (RQ4) se recurrió a pruebas pareadas —la $t$ de Student
para datos emparejados y la prueba de Wilcoxon, su equivalente no paramétrico—, que comparan
cada empresa consigo misma en ambos ejercicios y resultan, por ello, especialmente potentes
para detectar cambios temporales. Finalmente, para analizar los determinantes del discurso
(RQ3) se estimaron dos regresiones por mínimos cuadrados ordinarios con errores estándar
robustos a la heterocedasticidad: la primera explica el \gwindex{} en función del tamaño, los
ratios financieros, el sector, el año y la región; la segunda explica el tono financiero en
función de la especificidad climática y los mismos controles, con el fin de contrastar
directamente la relación entre optimismo y concreción. En ambos casos se verificó la ausencia
de problemas graves de multicolinealidad mediante el factor de inflación de la varianza. Los
resultados de todos estos contrastes se presentan y discuten en el capítulo siguiente.
